\documentclass[../main.tex]{subfiles}

\begin{document}

\section{Solutions for Note 4}
\begin{enumerate}
\item
For transformation $T( \langle a, b \rangle ) = a + b$, consider
the following equations for some scalars $\alpha$ and $\beta$.
\begin{align*}
    T ( \alpha \langle a_1, b_1 \rangle +
        \beta \langle a_2, b_2 \rangle )
    &= T ( \langle \alpha a_1, \alpha b_1 \rangle +
        \langle \beta a_2, \beta b_2 \rangle ) \\
    &= \alpha a_1 + \alpha b_1 + \beta a_2 + \beta b_2 \\
    &= \alpha (a_1 + b_1) + \beta (a_2 + b_2) \\
    &= \alpha T ( \langle a_1, b_1 \rangle ) +
        \beta T ( \langle a_2, b_2 \rangle )
\end{align*}
Therefore by the definition of linear transformation, $T$ is linear.

For $T( \langle a, b \rangle ) = ab$, consider the following for
arbitrary scalars $\alpha$ and $\beta$.
\begin{align*}
    T ( \alpha \langle a_1, b_1 \rangle +
        \beta \langle a_2, b_2 \rangle )
    &= T ( \langle \alpha a_1, \alpha b_1 \rangle +
        \langle \beta a_2, \beta b_2 \rangle ) \\
    &= T ( \langle \alpha a_1 + \beta b_2,
        \alpha b_1 + \beta b_2 \rangle ) \\
    &= (\alpha a_1 + \beta b_2)(\alpha b_1 + \beta b_2)
\end{align*}
Moreover,
\[
\alpha T ( \langle a_1, b_1 \rangle ) +
    \beta T ( \langle a_2, b_2 \rangle )
= \alpha a_1 b_1 + \beta a_2 b_2
\]
holds. Because $(\alpha a_1 + \beta b_2)(\alpha b_1 + \beta b_2) \neq
\alpha a_1 b_1 + \beta a_2 b_2$ in general, the transformation is not
linear.

\item
By definition, if $\lambda$ is an eigenvalue of $AB$, then there
exists some nonzero vector $v$ such that $AB v = \lambda v$.
First, let $\lambda \neq 0$. Notice that $u = Bv \neq \mathbf{0}$.
Therefore, $BA u = BA (Bv) = B (ABv) = B (\lambda v) = \lambda (Bv)
= \lambda u$ and $\lambda$ is also an eigenvalue of $BA$. Continuing
if $\lambda = 0$, then $\det(AB) = \det(A) \det(B) = 0 =
\det(B) \det(A) = \det(BA)$. In other words, if $AB$ is singular
then $BA$ is also singular and contains the eigenvalue $\lambda$.
Therefore, $\lambda$ is an eigenvalue of $AB$ if and only if it
is an eigenvalue of $BA$.

\item
First, to find bases for the kernel of the matrix, the following
equation must be solved.
\[
\begin{bmatrix}
    1 & 1 & 1 \\
    0 & 1 & 0
\end{bmatrix}
\begin{bmatrix}
    x \\
    y \\
    z
\end{bmatrix}
=
\begin{bmatrix}
    0 \\
    0
\end{bmatrix}
\]
Solving the system, $y = 0$ and $x = -z$. Therefore the basis for
the kernel is 
$\left\{
    \begin{bmatrix}
        -1 \\
        0 \\
        1
    \end{bmatrix}
\right\}$. Continuing, the basis for the image can be found.
Consider the following equation.
\[
\begin{bmatrix}
    1 & 1 & 1 \\
    0 & 1 & 0
\end{bmatrix}
\begin{bmatrix}
    x \\
    y \\
    z
\end{bmatrix}
=
\begin{bmatrix}
    x + y + z \\
    y
\end{bmatrix}
=
x
\begin{bmatrix}
    1 \\
    0
\end{bmatrix}
+ y
\begin{bmatrix}
    1 \\
    1
\end{bmatrix}
+ z
\begin{bmatrix}
    1 \\
    0
\end{bmatrix}
\]
Therefore,
\[
\left\{
\begin{bmatrix}
    1 \\
    0
\end{bmatrix},
\begin{bmatrix}
    1 \\
    1
\end{bmatrix}
\right\}
\]
is a basis of the image of the given matrix.

\item
By definition, $A \cong B$ if and only if there exists an invertible
matrix $P$ such that $A = P^{-1}BP$. Continuing, $A = P^{-1}BP$
if and only if $A^\intercal = (P^{-1}BP)^\intercal = P^\intercal
B^\intercal (P^{-1})^\intercal = P^\intercal B^\intercal
(P^\intercal)^{-1}$. In other words, $A \cong B$ if and only if
$A^\intercal \cong B^\intercal$.

\item
To find the eigenvalues for the given matrix say $A$, it suffices to
solve $\det(A - \lambda I) = 0$ where $\lambda$ is an eigenvalue.
\[
\det\left(
    \begin{bmatrix}
        1 - \lambda & 2           & 3 \\
        0           & 1 - \lambda & 0 \\
        0           & 1           & 2 - \lambda
    \end{bmatrix}
\right)
= 0
\]
Continuing with the expansion along the first column,
\begin{align*}
    \det\left(
        \begin{bmatrix}
            1 - \lambda & 2           & 3 \\
            0           & 1 - \lambda & 0 \\
            0           & 1           & 2 - \lambda
        \end{bmatrix}
    \right)
    &= (1 - \lambda)
    \det\left(
        \begin{bmatrix}
            1 - \lambda & 0 \\
            1           & 2 - \lambda
        \end{bmatrix}
    \right) \\
    &= (1 - \lambda)(1 - \lambda)(2 - \lambda)
\end{align*}
Therefore, the eigenvalues are $1$ and $2$. For eigenvalue $1$,
it suffices to solve the equation $Av = v$ or $(A - I)v = 0$.
Consider the following system.
\[
\begin{bmatrix}
    0 & 2 & 3 \\
    0 & 0 & 0 \\
    0 & 1 & 1
\end{bmatrix}
\begin{bmatrix}
    x \\
    y \\
    z
\end{bmatrix}
=
\begin{bmatrix}
    0 \\
    0 \\
    0
\end{bmatrix}
\]
Therefore, $2y + 3z = 0$ and $y + z = 0$. In other words, $y = z = 0$.
Thus, $S_1 =
\Span\left\{
    \begin{bmatrix}
        1 \\
        0 \\
        0
    \end{bmatrix}
\right\}$ holds for $\lambda = 1$. Continuing with $\lambda = 2$,
it suffices to solve $Av = 2v$ or $(A - 2I) v = 0$.
\[
\begin{bmatrix}
    -1 & 2  & 3 \\
    0  & -1 & 0 \\
    0  & 1  & 0
\end{bmatrix}
\begin{bmatrix}
    x \\
    y \\
    z
\end{bmatrix}
=
\begin{bmatrix}
    0 \\
    0 \\
    0
\end{bmatrix}
\]
Solving the system, $y = 0$ and $x = 3z$. Thus, $S_2 =
\Span\left\{
    \begin{bmatrix}
        3 \\
        0 \\
        1
    \end{bmatrix}
\right\}$ holds for $\lambda = 2$.
\end{enumerate}
\end{document}


